\begin{abstract}
The Routing Protocol for Low-Power and Lossy Networks (RPL) is the standard routing layer of industrial Internet of Things meshes, yet its lightweight defenses each address one control message at a time. A compromised node can flood DODAG Information Solicitations (DIS) while advertising a false parent in its Destination Advertisement Objects (DAO), and a DIS-only defense then detects the attacker but leaves downward routes broken. This paper proposes Adaptive Evidence Fusion for RPL (AEF-RPL), which extends the DIS-FDM Dempster--Shafer detector with a first-hop DAO-consistency evidence source whose uncertainty is self-calibrated, a persistence rule, selective DAO repair, and a 768-byte tabular Q-learning decision layer on every node. We evaluated AEF-RPL in Contiki-NG/Cooja over 2720 runs of 30-node static and mobile networks against DIS-FDM, Li-MSD, Secure-RPL, and RPL-MRC, using paired seeds, Wilcoxon tests with Holm correction, and a proposed Attack Resilience Index (ARI). Under the combined attack, AEF-RPL raised downward delivery and ARI over DIS-FDM in all six attacker-fraction and mobility conditions (for seven attackers in static networks, 98.3\% versus 77.2\% delivery and ARI 0.92 versus 0.71; adjusted $p=.012$), with equal attacker detection and no significant increase in false positives. On-node learning cut detection latency by 20--73\% but did not reduce false isolation, which was inherited from parent-change evidence under mobility. An ablation showed that selective DAO repair produced most of the service gain. All results can be regenerated from the archived logs with one command.
\end{abstract}
\begin{IEEEkeywords}
RPL security, intrusion detection, DIS flooding, DAO manipulation, Dempster--Shafer evidence fusion, Q-learning, network resilience, reproducible simulation.
\end{IEEEkeywords}
